\subsection{Explainable Models for Time-Series Diagnosis}

Explainable machine-learning methods aim to reveal which input variables contribute to a model prediction. Representative post-hoc approaches include LIME, Integrated Gradients, and SHAP, which explain predictions through local surrogate models, accumulated gradients, or additive feature contributions, respectively \cite{ribeiro2016lime,sundararajan2017axiomatic,lundberg2017shap}. For sequential data, TimeSHAP extends Shapley-based attribution to recurrent time-series models and can quantify the relevance of features, events, and time steps \cite{bento2021timeshap}. These methods have improved the interpretability of otherwise opaque predictive models and provide useful tools for analyzing multivariate time-series decisions.

Among these approaches, SHAP is particularly suitable for structured tabular representations because its additive attributions can be associated with individual input features. For tree ensembles, TreeSHAP provides an efficient formulation for computing such contributions and supports both local and global interpretation \cite{lundberg2017shap,lundberg2020trees}. This property is useful for UAV telemetry, where a flight window is represented by multiple statistical descriptors derived from physical sensor and state variables.

However, most existing explainability methods operate at the level of the model input representation. In UAV telemetry analysis, an attribution assigned to a statistical descriptor does not directly indicate which physical telemetry channel, uORB topic, or functional subsystem is most relevant to the diagnosed fault. Bridging this gap requires organizing feature-level attributions into engineering-level evidence that can be traced from model inputs to the underlying telemetry structure. HS-AeroTS builds on TreeSHAP in this direction by aggregating temporal-statistical contributions into progressively higher-level telemetry representations for subsequent architecture-aware diagnosis.
